% !TEX root = ../main.tex
%% Sujet type bac n° 7 — sujet inédit OptiTAB (calculs vérifiés par programme : bac_s7.py)

\sujetBac{7}{Probabilités conditionnelles/4,Fonction logarithme et suite/6,Géométrie : distance et volume/5,Intégration : aire et valeur moyenne/5}

\begin{exoBac}{4}{Probabilités conditionnelles}
On dispose de deux urnes. L'urne $U_1$ contient $3$ boules rouges et $2$ boules noires ; l'urne $U_2$ contient $1$ boule rouge et $4$ boules
noires. On lance un dé équilibré à six faces : si l'on obtient $1$ ou $2$, on tire une boule au hasard dans l'urne $U_1$, sinon dans
l'urne $U_2$. On note $U$ l'événement « le tirage a lieu dans l'urne $U_1$ » et $R$ l'événement « la boule tirée est rouge ».
\begin{questions}
  \item Représenter la situation par un arbre pondéré.
  \item Montrer que $P(R)=\frac13$.
  \item La boule tirée est rouge. Quelle est la probabilité qu'elle provienne de l'urne $U_1$ ?
  \item Les événements $U$ et $R$ sont-ils indépendants ?
  \item On répète cinq fois l'expérience, en remettant à chaque fois la boule tirée dans son urne. On note $Y$ le nombre de boules
    rouges obtenues. Préciser la loi de $Y$, puis calculer la probabilité d'obtenir au moins une boule rouge et celle d'obtenir
    exactement deux boules rouges. On donnera les valeurs exactes puis arrondies au millième.
  \item Combien de fois au minimum faut-il répéter l'expérience pour que la probabilité d'obtenir au moins une boule rouge dépasse
    $0{,}999$ ?
\end{questions}
\end{exoBac}

\begin{exoBac}{6}{Fonction logarithme et suite}
\partieExo{A}{Étude d'une fonction}
On considère la fonction $f$ définie sur $\intOO0{+\infty}$ par $f(x)=x-\ln x$.
\begin{questions}
  \item Déterminer les limites de $f$ en $0$ et en $+\infty$.
  \item Étudier les variations de $f$ et dresser son tableau de variations.
  \item En déduire que, pour tout réel $x>0$, $\ln x\leq x-1$.
\end{questions}
\partieExo{B}{Une suite récurrente}
On considère la suite $(u_n)$ définie par $u_0=2$ et, pour tout entier naturel $n$, $u_{n+1}=u_n-\ln u_n$.
\begin{questions}[resume]
  \item Calculer $u_1$ et $u_2$ ; on en donnera des valeurs approchées au millième.
  \item Démontrer par récurrence que, pour tout entier naturel $n$, $1\leq u_{n+1}\leq u_n$.
  \item En déduire que la suite $(u_n)$ converge, puis déterminer sa limite.
  \item On considère la fonction Python ci-contre.
\begin{codecote}[0.44]
\begin{pythoncol}
from math import log

def rang(e):
    u = 2
    n = 0
    while u - 1 >= e:
        u = u - log(u)
        n = n + 1
    return n
\end{pythoncol}
\cote
\begin{enumerate}[label=\alph*),leftmargin=1.4em,itemsep=0.6mm,topsep=0mm]
  \item Que représente la valeur renvoyée par \code{rang(e)} ?
  \item Que renvoie \code{rang(1e-6)} ? On justifiera à l'aide des valeurs successives de $u_n$.
  \item Que penser de la vitesse de convergence de cette suite ?
\end{enumerate}
\end{codecote}
\end{questions}
\end{exoBac}

\begin{exoBac}{5}{Géométrie : distance et volume}
\begin{codecote}[0.6]
L'espace est muni d'un repère orthonormé d'origine $O$. On considère les points $A(2\pv0\pv0)$, $B(0\pv4\pv0)$, $C(0\pv0\pv4)$ et
$D(1\pv3\pv3)$.
\begin{questions}
  \item Justifier que les points $A$, $B$, $C$ définissent un plan, puis vérifier que ce plan a pour équation $2x+y+z-4=0$.
  \item Déterminer une représentation paramétrique de la droite $\Delta$ passant par $D$ et orthogonale au plan $(ABC)$.
\end{questions}
\cote
\centering
\begin{tikzpicture}[otEspace,scale=0.62]
  \coordinate (O) at (0,0,0); \coordinate (A) at (2,0,0); \coordinate (B) at (0,4,0); \coordinate (C) at (0,0,4);
  \coordinate (D) at (1,3,3);
  \draw[otMuted,line width=0.4pt,-{Stealth[length=1.6mm]}] (O) -- (3,0,0);
  \draw[otMuted,line width=0.4pt,-{Stealth[length=1.6mm]}] (O) -- (0,5,0);
  \draw[otMuted,line width=0.4pt,-{Stealth[length=1.6mm]}] (O) -- (0,0,5);
  \fill[otNavyTint,opacity=0.8] (A) -- (B) -- (C) -- cycle;
  \draw[otVu] (A) -- (B) -- (C) -- cycle;
  \draw[otGreenDark,line width=0.8pt] (D) -- (A) (D) -- (C);
  \draw[otGreenDark,line width=0.8pt,dash pattern=on 2.2pt off 1.6pt] (D) -- (B);
  \node[otPtV] at (D) {};
  \foreach \p/\a in {O/north east,A/north,B/west,C/east,D/south west}{\node[otSommet,anchor=\a] at (\p) {$\p$};}
\end{tikzpicture}
\end{codecote}
\begin{questions}[start=3]
  \item Déterminer les coordonnées du point $H$, intersection de $\Delta$ et du plan $(ABC)$, puis la distance du point $D$ au plan $(ABC)$.
  \item Calculer $\vect{AB}\cdot\vect{AC}$, puis $\cos\widehat{BAC}$ et $\sin\widehat{BAC}$. En déduire que l'aire du triangle $ABC$ vaut
    $4\sqrt6$.
  \item Calculer le volume du tétraèdre $ABCD$.
  \item Calculer la distance du point $O$ au plan $(ABC)$ et le volume du tétraèdre $OABC$. Expliquer pourquoi les tétraèdres $OABC$ et
    $DABC$ ont le même volume.
\end{questions}
\end{exoBac}

\begin{exoBac}{5}{Intégration : aire et valeur moyenne}
On note $\mathcal{C}$ la courbe représentative de la fonction exponentielle dans un repère orthonormé, et $T$ sa tangente au point
d'abscisse $0$.
\begin{questions}
  \item Déterminer une équation de $T$, puis démontrer que, pour tout réel $x$, $e^{x}\geq x+1$.
  \item Calculer l'aire, en unités d'aire, du domaine compris entre $\mathcal{C}$, la droite $T$ et les droites d'équations $x=0$ et $x=1$.
    On en donnera la valeur exacte, puis une valeur approchée au centième.
  \item Pour tout réel $a>0$, on note $A(a)$ l'aire du domaine compris entre $\mathcal{C}$, $T$ et les droites d'équations $x=0$ et $x=a$.
    Exprimer $A(a)$ en fonction de $a$, puis déterminer la limite de $A(a)$ quand $a$ tend vers $+\infty$.
  \item Montrer que l'équation $A(a)=1$ admet une unique solution sur $\intOO0{+\infty}$, et en donner un encadrement d'amplitude $10^{-2}$.
  \item Calculer la valeur moyenne de la fonction exponentielle sur l'intervalle $\intFF0{\ln4}$.
\end{questions}
\end{exoBac}
